\documentclass{beamer}
% COMSC-3013 Computer Architecture shared Beamer preamble.
% Week decks live one directory below weeks/ and input this file as:
%   % COMSC-3013 Computer Architecture shared Beamer preamble.
% Week decks live one directory below weeks/ and input this file as:
%   % COMSC-3013 Computer Architecture shared Beamer preamble.
% Week decks live one directory below weeks/ and input this file as:
%   \input{../_shared/beamer-preamble.tex}
% Keep the course self-contained. Change visual/build conventions here,
% not by forking one-off week themes.

\usetheme{Boadilla}
\usecolortheme{seahorse}
\setbeamertemplate{navigation symbols}{}
\setbeamertemplate{footline}[frame number]

\usepackage{amsmath}
\usepackage{booktabs}
\usepackage{graphicx}
\usepackage{listings}
\usepackage{pgfpages}

\lstset{
  basicstyle=\ttfamily\small,
  columns=fullflexible,
  keepspaces=true,
  showstringspaces=false,
  breaklines=true
}

% Speaker notes remain hidden in the ordinary student build. Define \shownotes
% before inputting the deck to create an instructor copy; see weeks/README.md.
\ifdefined\shownotes
  \setbeameroption{show notes on second screen=right}
\fi

\newcommand{\ArchTitleSlide}[4]{%
  % #1 week label, #2 title, #3 central question, #4 Wednesday prediction
  \begin{frame}
    \titlepage
  \end{frame}
  \begin{frame}{The machine question}
    \Large #3

    \vspace{1.2em}
    \normalsize\textbf{Prediction to carry into Wednesday:} #4
  \end{frame}
}

% Keep the course self-contained. Change visual/build conventions here,
% not by forking one-off week themes.

\usetheme{Boadilla}
\usecolortheme{seahorse}
\setbeamertemplate{navigation symbols}{}
\setbeamertemplate{footline}[frame number]

\usepackage{amsmath}
\usepackage{booktabs}
\usepackage{graphicx}
\usepackage{listings}
\usepackage{pgfpages}

\lstset{
  basicstyle=\ttfamily\small,
  columns=fullflexible,
  keepspaces=true,
  showstringspaces=false,
  breaklines=true
}

% Speaker notes remain hidden in the ordinary student build. Define \shownotes
% before inputting the deck to create an instructor copy; see weeks/README.md.
\ifdefined\shownotes
  \setbeameroption{show notes on second screen=right}
\fi

\newcommand{\ArchTitleSlide}[4]{%
  % #1 week label, #2 title, #3 central question, #4 Wednesday prediction
  \begin{frame}
    \titlepage
  \end{frame}
  \begin{frame}{The machine question}
    \Large #3

    \vspace{1.2em}
    \normalsize\textbf{Prediction to carry into Wednesday:} #4
  \end{frame}
}

% Keep the course self-contained. Change visual/build conventions here,
% not by forking one-off week themes.

\usetheme{Boadilla}
\usecolortheme{seahorse}
\setbeamertemplate{navigation symbols}{}
\setbeamertemplate{footline}[frame number]

\usepackage{amsmath}
\usepackage{booktabs}
\usepackage{graphicx}
\usepackage{listings}
\usepackage{pgfpages}

\lstset{
  basicstyle=\ttfamily\small,
  columns=fullflexible,
  keepspaces=true,
  showstringspaces=false,
  breaklines=true
}

% Speaker notes remain hidden in the ordinary student build. Define \shownotes
% before inputting the deck to create an instructor copy; see weeks/README.md.
\ifdefined\shownotes
  \setbeameroption{show notes on second screen=right}
\fi

\newcommand{\ArchTitleSlide}[4]{%
  % #1 week label, #2 title, #3 central question, #4 Wednesday prediction
  \begin{frame}
    \titlepage
  \end{frame}
  \begin{frame}{The machine question}
    \Large #3

    \vspace{1.2em}
    \normalsize\textbf{Prediction to carry into Wednesday:} #4
  \end{frame}
}


\title{Week 5 Monday}
\subtitle{Build the Machine}
\author{COMSC-3013 Computer Architecture}
\date{}

\begin{document}

\begin{frame}
  \titlepage
\end{frame}

\ArchQuestionSlide{What should I build for this workload, what does each part buy me, and how should I compare choices?}{Under a fixed \$1,500 tower budget, which resource deserves the largest budget share for your workload, and what evidence would make you move about \$150?}
\note{Open with the student's useful prior belief: bigger numbers often buy capability. The week is not here to shame that instinct; it is here to make the word ``better'' conditional on workload and evidence.}

\begin{frame}{Start with verbs, not parts}
\begin{itemize}
  \item Compile, render, encode, analyze, store, move, respond.
  \item A workload names what the machine must \textbf{do}.
  \item A specification describes a resource a component provides.
  \item The design question is whether that resource matters under this workload and budget.
\end{itemize}
\note{Use two different workloads and ask whether the same expensive component deserves the same share of the budget. Do not begin naming a preferred CPU/GPU.}
\end{frame}

\begin{frame}{First gate: would it run?}
\small
\begin{tabular}{ll}
\toprule
Relationship & Bounded compatibility question \\
\midrule
CPU $\leftrightarrow$ board & socket/platform/firmware support? \\
Board $\leftrightarrow$ RAM & correct memory type/capacity? \\
Storage $\leftrightarrow$ board & compatible slot/interface? \\
Board/GPU/cooler $\leftrightarrow$ case & form factor and clearance? \\
CPU/GPU $\leftrightarrow$ PSU & power budget and connectors? \\
Graphics path & integrated or discrete path exists? \\
\bottomrule
\end{tabular}
\note{A commercial builder can help surface problems, but the manufacturer specs explain the actual relationship. Keep this bounded. We are not teaching every physical/electrical PC-building detail.}
\end{frame}

\begin{frame}{Three questions that are not the same}
\begin{columns}[T]
\column{0.31\textwidth}
\textbf{Capacity}

How much can be held?

\column{0.31\textwidth}
\textbf{Latency}

How long until one requested thing is available?

\column{0.31\textwidth}
\textbf{Bandwidth}

How much can move/complete per unit time once work is flowing?
\end{columns}

\vspace{1.4em}
\centering
A bigger number on one axis does not answer the other two.
\note{These words return all semester. Do not teach numeric latency tables here. The goal is simply to prevent students from treating GB, ns, and GB/s as versions of the same idea.}
\end{frame}

\begin{frame}{The hierarchy starts a story}
\small
\begin{tabular}{llll}
\toprule
Layer & Capacity tendency & Buy separately? & Persistent? \\
\midrule
L1/L2/L3 cache & small & no & no \\
RAM & much larger & yes & no \\
SSD/NVMe & larger again & yes & yes \\
optional bulk storage & very large & yes & yes \\
\bottomrule
\end{tabular}

\vspace{1em}
\textbf{Week 5 question:} what can we responsibly compare now?\\
\textbf{Week 10 question:} what happens when we make the hierarchy hurt?
\note{This is explicitly a first map. Cache has no honest retail price-per-GB. Week 10 will add measured sensitivity instead of memorized internet numbers.}
\end{frame}

\begin{frame}{Dollars Per: correct arithmetic, weak verdict}
\begin{center}
\begin{tabular}{lrrr}
\toprule
Synthetic option & Price & Cores & Naive \$/core \\
\midrule
CPU A & \$280 & 8 & \$35.00 \\
CPU B & \$390 & 12 & \$32.50 \\
\bottomrule
\end{tabular}
\end{center}

CPU B wins this metric. What has the metric \textbf{not} told us?
\note{Stress that the numbers are invented. Ask for missing variables: workload scaling, core type, single-thread behavior, memory behavior, power, and opportunity cost. The arithmetic is true. The implied judgment is not yet earned.}
\end{frame}

\begin{frame}{Some Dollars-Per is literal}
\begin{itemize}
  \item RAM \$/GB: real separately purchased capacity.
  \item Storage \$/TB: real separately purchased capacity.
  \item CPU \$/core or GPU \$/GB VRAM: intentionally naive workload proxies.
  \item L1/L2/L3 cache \$/GB: \textbf{do not invent it}.
\end{itemize}

\vspace{1em}
Cache is an architectural scarcity problem: capacity competes with area, latency, power, organization, and total chip design.
\note{This is the correction to the historical Dollars-Per activity. We preserve the instinct but stop pretending cache is a retail SSD glued to the CPU.}
\end{frame}

\begin{frame}{A budget is an architecture statement}
\[
\text{budget share} = \frac{\text{component price}}{1500}\times 100\%
\]

\begin{itemize}
  \item 40\% on one component is a claim about what the workload needs.
  \item Every dollar has an opportunity cost.
  \item The same parts can be sensible for one workload and silly for another.
\end{itemize}
\note{The fixed budget equalizes the intellectual task. Owning a more expensive real computer confers no advantage.}
\end{frame}

\begin{frame}{Make tradeoff a verb}
Build one compatible baseline, then move about \$100--\$200 while keeping workload and total budget approximately fixed.

\vspace{1em}
Examples:
\begin{itemize}
  \item CPU $\rightarrow$ GPU;
  \item storage capacity $\rightarrow$ RAM;
  \item discrete GPU $\rightarrow$ CPU/RAM/storage when integrated graphics is enough;
  \item premium case/board $\rightarrow$ workload-facing resource.
\end{itemize}
\note{Students are not asked for two giant builds. One bounded reallocation creates the comparison. Ask: what became better, what became worse, and what evidence is still missing?}
\end{frame}

\begin{frame}{AI Fluency Lens 5: select the right model}
\begin{itemize}
  \item Ask for recommendations under meaningfully different workload framing.
  \item Identify the assumed workload and metric behind each recommendation.
  \item Verify exact compatibility/spec claims outside the model.
  \item Treat recommendation confidence as rhetoric, not evidence.
\end{itemize}

\vspace{1em}
\centering
\textbf{AI can propose. Evidence decides.}
\note{This can use two providers, one provider with two prompts, or course-provided samples. No paid AI dependency. Verify one claim live against an official manufacturer specification page.}
\end{frame}

\begin{frame}{What Week 5 can and cannot claim}
\textbf{Can support:}
\begin{itemize}
  \item compatibility under checked evidence;
  \item dated market price;
  \item capacity and budget-share comparisons;
  \item a workload hypothesis.
\end{itemize}

\textbf{Cannot yet support:}
\begin{itemize}
  \item measured application performance;
  \item cache latency/bandwidth;
  \item scaling or accelerator benefit;
  \item ``best PC'' conclusions.
\end{itemize}
\note{This is the epistemic boundary. The rest of the technical arc exists because Week 5 does not have enough evidence yet.}
\end{frame}

\begin{frame}{Wednesday: make the design move}
\Large
Build the \$1,500 baseline. Move about \$150. Keep the workload fixed.

\vspace{1em}
\normalsize
Record what the move buys, what it costs, and what you still cannot prove.

\vspace{1em}
\textbf{Prediction to preserve:} Which resource deserves the budget, and what evidence would make you reallocate it?
\note{Stop here. Do not show a canonical ``correct'' parts list. Wednesday's purpose is for the student to own the first hypothesis.}
\end{frame}

\end{document}
